% Copyright (c) 2026 Twin Vector Labs LLC.
% All rights reserved.
%
% Research-direction companion to spin_readout.tex.  Takes seriously the idea of
% making the LOOPS that encircle pairs of color holes the primary objects ("the
% quarks") instead of the holes.  Three results: (i) a pair-encircling cycle is
% the Poincare dual of the complementary hole (homological additivity + the
% singlet constraint), so "loops as quarks" is the EDGE basis of the same three
% quarks, not a new set; (ii) flavor read on the loops gives uud with the
% diquark-loop carrying d -- the Poincare dual of the d-hole, matching the
% two-step diquark build (#489); (iii) a loop around two QUARK holes is a diquark
% (3-bar), not a gluon -- a gluon is a quark-antiquark loop, which exists across
% the proton/antiproton legs of the pair-creation build.  This is a lens, not a
% validated result; the spin readout it motivates REDUCES to the closed-loop
% holonomy and pairwise C_ij of cartan_weyl_gluon.tex (which absorbed the retired
% total-space-operator note) -- it does not bypass its hard fiber<>cells kernel.
\documentclass[11pt,oneside]{article}

\usepackage[margin=1in]{geometry}
\usepackage{amsmath}
\usepackage{amssymb}
\usepackage{amsthm}
\usepackage{mathtools}
\usepackage{braket}
\usepackage{bm}
\usepackage{enumitem}
\usepackage{booktabs}
\usepackage{array}
\usepackage{xcolor}
\usepackage[colorlinks=true,linkcolor=blue!55!black,citecolor=green!45!black,%
            urlcolor=blue!55!black]{hyperref}

% --- convenience macros (mirrors spin_readout.tex) ---------------------------
\newcommand{\R}{\mathbb{R}}
\newcommand{\C}{\mathbb{C}}
\newcommand{\Z}{\mathbb{Z}}
\newcommand{\halff}{\tfrac{1}{2}}
\newcommand{\id}{\mathbb{1}}
\newcommand{\Sig}{\Sigma}
\newcommand{\Tr}{\operatorname{Tr}}
\newcommand{\eqdef}{\vcentcolon=}
\DeclareMathOperator{\sgn}{sgn}
\newcommand{\expval}[1]{\langle #1 \rangle}
\newcommand{\mel}[3]{\langle #1 | #2 | #3 \rangle}
\newcommand{\norm}[1]{\left\lVert #1 \right\rVert}
\newcommand{\code}[1]{\texttt{#1}}
\newcommand{\up}{\uparrow}
\newcommand{\down}{\downarrow}
% additions for this companion
\newcommand{\SU}{\mathrm{SU}}
\newcommand{\SO}{\mathrm{SO}}
\newcommand{\Spin}{\mathrm{Spin}}
\newcommand{\bn}{\bar{\mathbf{3}}}
\newcommand{\thr}{\mathbf{3}}
\newcommand{\eg}{\mathbf{8}}
\newcommand{\six}{\mathbf{6}}
\theoremstyle{plain}
\newtheorem{claim}{Claim}
\newtheorem{remark}{Remark}

\title{\bfseries Pair--Encircling Loops as Quarks\\[4pt]
\large The Poincar\'e--dual edge basis, flavor from the diquark loop,\\
and why a two--quark loop is a diquark, not a gluon}
\author{Tessera theory notes \quad(companion to \code{spin\_readout.tex}; part of issue \#410)}
\date{\today}

\begin{document}
\maketitle

\begin{abstract}
\noindent
This is a research--direction companion to \code{spin\_readout.tex} (same
directory). It takes seriously a single proposal: make the \emph{loops that
encircle pairs} of color holes the primary objects --- ``the quarks'' --- rather
than the holes themselves, and ask what flavor and what gauge content they carry.
Three things follow, two of them theorems and one a reframing. (i) A
pair--encircling cycle is, by homological additivity together with the carried
singlet constraint $w_A+w_B+w_C=0$, the Poincar\'e \emph{dual} of the
\emph{complementary} hole: $[\,\text{loop around }i,j\,]=-[\,\text{hole }k\,]$. So
``loops as quarks'' is not a new set of three objects; it is the \emph{edge
basis} of the very same three quarks, carrying the identical $[1,\omega,\omega^2]$
color singlet. (ii) Flavor read on the loops gives $uud$ with the
\emph{diquark loop} carrying $d$ --- and this is exactly the dual of the
$d$--hole (the quark \emph{outside} the diquark), so it matches the two--step
diquark build of \#489 with no inversion and no fiat; the loop combinatorics
reproduce the proton's mixed--symmetry skeleton under the diquark's internal
exchange. (iii) Group theory settles the ``gluon'' alternative: a loop around two
\emph{quark} holes lives in $\thr\otimes\thr=\six\oplus\bn$ --- a \emph{diquark}
($\bn$), \emph{not} a gluon ($\thr\otimes\bn=\eg\oplus\mathbf 1$); a genuine
gluon loop must enclose a quark \emph{and an antiquark}, which exist across the
proton/antiproton legs of the pair--creation build. \textbf{Status: a lens, not a
validated result.} The \emph{spin} readout this picture motivates reduces to the
closed--loop holonomy and pairwise correlator of \code{cartan\_weyl\_gluon.tex};
it sharpens \emph{what} the loops
are and \emph{which} flavor/color they carry, and it does \emph{not} bypass the
hard fiber$\leftrightarrow$cells kernel those notes name.
\end{abstract}

\tableofcontents
\bigskip

% =============================================================================
\section{Scope: what this companion adds, and what it does not}
\label{sec:scope}
% =============================================================================

Three sibling documents in this directory have already settled the surrounding
problem, and we build directly on them rather than restating them. We write
\textbf{[SR\,\S$n$]} and \textbf{[CW\,\S$n$]} for sections of,
respectively, \code{spin\_readout.tex} (the study document: the proton's
$J^2=\tfrac34$ is an entangled mixed--symmetry quantity, and a per--hole product
read provably floors at the mixture $\sim\tfrac94$; it names the cure, the
total--space operator $\hat J^2=\tfrac34+\sum_{i<j}\hat P_{ij}$ keeping the
cross--hole terms) and \code{cartan\_weyl\_gluon.tex} (the Cartan/Weyl lens: the
color group's Weyl $S_3$ \emph{is} the swap $S_3$; the missing object is the
abelian $A$ factor, coordinatized by the connected pairwise correlator $C_{ij}$;
the entangling exponential acts like the one--gluon--exchange hyperfine
interaction; and the fiber$\leftrightarrow$cells lift with the closed--loop
holonomy ``escape hatch'' --- the consolidated home of the retired
total--space note).

\paragraph{The one new move.}
All three siblings keep the three \emph{holes} as the objects and add pairwise
\emph{operators} (swaps $P_{ij}$, correlators $C_{ij}$) on top. This note instead
takes the user's proposal literally: promote the three \emph{loops}
$\{\gamma_{AB},\gamma_{AC},\gamma_{BC}\}$ encircling the three pairs to the
primary degrees of freedom and ask what they are. The payoff is not a new spin
mechanism --- \S\ref{sec:spin} shows the spin readout reduces to [CW\,\S6] ---
but a sharper account of (a) the \emph{topology} of the reframe
(\S\ref{sec:duality}: it is the Poincar\'e--dual edge basis), (b) \emph{flavor}
(\S\ref{sec:flavor}: the diquark loop is $d$, dual to the $d$--hole), and (c) the
\emph{gluon} question (\S\ref{sec:gluon}: rep theory makes the loop a diquark).

\paragraph{Status discipline.}
\S\ref{sec:duality} and \S\ref{sec:gluon} are theorems (homological additivity;
$\thr\otimes\thr$ Clebsch--Gordan). \S\ref{sec:flavor} is a consistent
\emph{labeling} that reproduces the proton's symmetry skeleton --- necessary, not
sufficient, for $\tfrac34$. The reframe does not, by itself, produce a number;
\S\ref{sec:caveats} states exactly where it reduces to the already--open kernel.

% =============================================================================
\section{The pair--encircling cycle is the dual of the complementary hole}
\label{sec:duality}
% =============================================================================

\paragraph{Setup.}
The host is a closed combinatorial $S^4$ with three color holes $A,B,C$ (three
disjoint removed top cells, $b_3=n-1=2$ independent classes for $n=3$ holes)
[SR\,\S5--7]. The carried representative is a single harmonic $3$--cochain
$\psi=\code{EigenstateSynthesis(st,3).carriedRepresentative}$, and each hole $h$
is detected by a cycle $[h]$ (its link, a separating $3$--sphere; on the dual
$1$--skeleton, the closed BFS path \code{wilson\_line} already walks) whose period
$w_h\eqdef\oint_{[h]}\psi$ is the carried weight. The color singlet is
\begin{equation}
  (w_A,w_B,w_C)=(1,\omega,\omega^2),\qquad \omega=e^{2\pi i/3},
  \qquad w_A+w_B+w_C=0,
  \label{eq:singlet}
\end{equation}
the discrete $\epsilon_{abc}$ [SR\,\S4,\S6]; the singlet $[1,1,1]$ is \emph{not}
carried (residual $\sim100$), which is confinement.

\paragraph{Additivity.}
A cycle $\gamma_{ij}$ enclosing \emph{both} holes $i$ and $j$ bounds a region with
those two punctures, hence is homologous to the sum of the two link cycles, and
periods are linear over cycles:
\begin{equation}
  [\gamma_{ij}] \;=\; [i]+[j],
  \qquad
  \oint_{\gamma_{ij}}\psi \;=\; w_i+w_j .
  \label{eq:additivity}
\end{equation}
This is the incidence (boundary) map of the triangle $K_3$ on $\{A,B,C\}$ sending
each \emph{edge} to the sum of its two endpoint \emph{nodes}.

\paragraph{The duality.}
Impose the singlet \eqref{eq:singlet} in \eqref{eq:additivity}:
\begin{equation}
  \boxed{\;[\gamma_{AB}]=-[C],\quad [\gamma_{AC}]=-[B],\quad [\gamma_{BC}]=-[A].\;}
  \label{eq:dual}
\end{equation}
\emph{The loop around a pair is minus the complementary hole.} Concretely the
three pair--loops carry $(w_{AB},w_{AC},w_{BC})=(-\omega^2,-\omega,-1)$, which is
the \emph{same} register $[1,\omega,\omega^2]$ up to an overall sign and a
relabeling. Two immediate consequences:
\begin{itemize}[leftmargin=1.6em]
  \item \textbf{Color is unchanged.} The loop basis carries the identical singlet,
        dual--labeled. So for \emph{color} the reframe is a relabeling and adds
        nothing --- a consistency check, not new physics. (Good: any framework in
        which ``loops as quarks'' broke the carried singlet would be wrong.)
  \item \textbf{No new degrees of freedom.} With $\sum_h[h]=0$ from the singlet
        and the dual relation $\sum_{i<j}[\gamma_{ij}]=2\sum_h[h]=0$, the three
        pair--loops span the \emph{same} two--dimensional space as the three holes
        (mod the singlet), with one linear relation each. As $S_3$
        representations both the node space and the edge space are
        $\mathbf 1\oplus\mathbf 2$ (trivial $\oplus$ standard): identical content,
        different basis.
\end{itemize}

\begin{remark}[The minus sign is the antitriplet]
The conjugation $w_{ij}=-w_k$ in \eqref{eq:dual} is not cosmetic: it is the
statement that the pair transforms as the \emph{conjugate} of the complementary
quark. This is exactly the $\bn$ (antitriplet) diquark of \S\ref{sec:gluon}, seen
here homologically before any group theory.
\end{remark}

% =============================================================================
\section{Loops as quarks is the edge basis --- and the spin payoff reduces to
the siblings}
\label{sec:spin}
% =============================================================================

By \S\ref{sec:duality} ``treat the loops as the quarks'' means: change from the
\emph{node} basis $\{A,B,C\}$ to the \emph{edge} basis
$\{\gamma_{AB},\gamma_{BC},\gamma_{CA}\}$ of the same register. Why bother, if it
adds no degrees of freedom? Because the natural \emph{observables} differ:
\begin{itemize}[leftmargin=1.6em]
  \item the node basis invites \emph{per--hole} (diagonal) reads --- one spinor per
        hole, \code{kron}'d --- which is exactly \code{composite\_j2} and
        \emph{provably} the product floor [SR\,\S8];
  \item the edge basis invites \emph{pairwise} reads, and $J^2=\tfrac34+\sum_{i<j}
        P_{ij}$ \emph{is} a sum over pairs [SR\,\S3]. The edge basis is the one in
        which the off--diagonal swap structure of $J^2$ is diagonal.
\end{itemize}
This is the right instinct, but it is the \emph{same} instinct the siblings
already act on, and the honest accounting is that the \emph{spin} content of a
pair--loop reduces to objects they define. There are two distinct loop
observables, and only one carries spin:
\begin{enumerate}[leftmargin=1.9em]
  \item \textbf{The period} $\oint_{\gamma_{ij}}\psi=w_i+w_j$ (\eqref{eq:additivity}).
        This is \emph{linear} in $\psi$ --- a one--body datum in the internal
        color/taste space. It gives the duality \eqref{eq:dual} and the flavor of
        \S\ref{sec:flavor}, but \emph{not} the spin correlation, which is bilinear.
  \item \textbf{The closed--loop holonomy} --- the $\Spin(4)$ rotation accumulated
        transporting a spinor around $\gamma_{ij}$. This \emph{is} the bilocal,
        cross--hole object. But it is precisely the closed--loop holonomy
        [CW\,\S6] names as the entangling carrier (its ``escape hatch''): the
        pair--encircling loop $\gamma_{ij}$ is the exchange loop
        $\Gamma_{ij}\cup\Gamma_{ji}$ of the swap $\hat P_{ij}$ in the swap form
        $J^2=\tfrac34+\sum_{i<j}P_{ij}$ [CW\,\S3, Eq.\,(j2swap)], and its spinor
        lift is the hinge holonomy $W_{\text{spin}}(h)=\cos(\varepsilon_h/2)$ of
        \code{dk\_spin\_wilson.py} (the \#477 spin--holonomy thread).
\end{enumerate}

\begin{claim}[The spin reframe is a renaming of {[CW]}]
Reading a spinor on the pair--loop, done correctly (the holonomy, not the period),
is the closed--loop--holonomy swap $\hat P_{ij}$ of [CW\,\S6] (the escape hatch);
summing the three gives $\hat J^2=\tfrac34+\sum_{i<j}\hat P_{ij}$. Equivalently, the connected
pairwise correlator $C_{ij}=\expval{\bm S_i\!\cdot\!\bm S_j}-\expval{\bm S_i}\!\cdot\!
\expval{\bm S_j}$ of [CW\,\S5] is the chamber coordinate of the same loop. The
edge basis gives these a \emph{geometric name} (``the loop around the pair''); it
does not give a new mechanism.
\end{claim}

So the value of ``loops as quarks'' is \emph{not} on the spin side, where it
reduces to known plans. It is on the \emph{flavor} and \emph{color} sides, which
the siblings treat only glancingly, and which the edge basis settles cleanly. The
rest of this note is those two.

% =============================================================================
\section{Flavor from the diquark loop --- and why it is not an inversion}
\label{sec:flavor}
% =============================================================================

\paragraph{The construction.}
The proton is \emph{built}, not posited [SR\,\S7], in two steps [\#489]:
$q+q\to\text{diquark}$, then $\text{diquark}+q\to\text{proton}$. Let holes $1,2$
be the two quarks that bond in step~1 (the diquark constituents) and hole $3$ the
quark added in step~2. The user's proposal assigns flavor to the three
\emph{loops}:
\begin{equation}
  \gamma_{12}\ (\text{the diquark loop})\;=\;\ket d,\qquad
  \gamma_{13}=\ket u,\qquad \gamma_{23}=\ket u .
  \label{eq:loopflavor}
\end{equation}
The multiplicity is $uud$ --- one $d$, two $u$ --- the proton's. The question is
whether this is consistent or a sleight of hand, since in the standard picture the
\emph{diquark} is the like--flavor $uu$ pair, so calling the diquark loop ``$d$''
looks backwards.

\paragraph{It is the Poincar\'e dual of the node flavor, and it agrees.}
It is not backwards. By the duality \eqref{eq:dual}, $\gamma_{12}=-[3]$: the
diquark loop \emph{is} the dual of hole $3$, the quark \emph{outside} the diquark.
In the standard $uud$ build the diquark is $\{1,2\}=uu$ and the outside quark is
$3=d$. Hence
\begin{equation}
  \gamma_{12}=-[3]=-[\,d\text{--hole}\,],\quad
  \gamma_{13}=-[2]=-[\,u\text{--hole}\,],\quad
  \gamma_{23}=-[1]=-[\,u\text{--hole}\,],
\end{equation}
so the loop flavor $\{\gamma_{12},\gamma_{13},\gamma_{23}\}=\{d,u,u\}$ is exactly
the Poincar\'e dual of the node flavor $\{1,2,3\}=\{u,u,d\}$. \textbf{The diquark
loop carries $d$ precisely \emph{because} it is dual to the $d$--quark.} No
inversion, no fiat --- and, per the codebase discipline, flavor is \emph{read off}
the dual basis, not imposed [\,cf.\ the removal of \code{relax\_singlet} for
hand--placing the singlet, CW\,\S6 caveat\,].

\paragraph{The mixed--symmetry skeleton, for free.}
Let $\tau=(1\!\leftrightarrow\!2)$ be the exchange of the two diquark
constituents. On the loops,
\begin{equation}
  \tau:\quad \gamma_{12}\mapsto\gamma_{12}\ (\text{fixed}),\qquad
  \gamma_{13}\leftrightarrow\gamma_{23}\ (\text{swapped}).
  \label{eq:tau}
\end{equation}
The $d$--loop is the $\tau$--spectator; the two $u$--loops are the $\tau$--exchanged
pair. This is exactly the proton's defining symmetry --- the mixed--symmetry state
$2\ket{uud}-\ket{udu}-\ket{duu}$ is symmetric under $1\!\leftrightarrow\!2$ with
the $d$ in the distinguished slot [SR\,\S3--4, Eq.\,(su6)]. The edge basis
\emph{produces the proton's combinatorial skeleton} from the bare topology of
which two holes bonded first. Contrast the node read, whose blind per--hole
symmetrization lands on the fully symmetric quartet or the product mixture
[SR\,\S4,\S8]: the loop labeling supplies precisely the structure the product read
destroys.

\paragraph{Honest limit.}
This is the right \emph{skeleton}, which is \emph{necessary} for $\tfrac34$; it is
not \emph{sufficient}. The actual value depends on the pairwise spin correlations
--- the bilinear holonomy/$C_{ij}$ data of \S\ref{sec:spin}, which the period
\eqref{eq:loopflavor} does not contain. The loop flavor fixes the
\emph{representation type} ($[2,1]$ mixed, the proton channel rather than the
$[3]$ symmetric $\Delta$ channel); recovering the number $\tfrac34$ still requires
[CW]'s closed--loop holonomy. The contribution here is that the diquark
build [\#489] and the proton's spin--flavor symmetry are \emph{the same fact} seen
in two bases.

% =============================================================================
\section{Loops as gluons? Rep theory says diquark}
\label{sec:gluon}
% =============================================================================

The user's alternative --- ``what if the loops were gluons?'' --- is settled by
the color representation of what a loop encloses. Two facts of $\SU(3)$
Clebsch--Gordan decide it.

\paragraph{A two--quark loop is a diquark ($\bn$), not a gluon.}
By Gauss's law the color flux through a surface bounding $\gamma_{ij}$ equals the
total enclosed charge, here two quarks:
\begin{equation}
  \thr\otimes\thr \;=\; \bn \oplus \six
  \qquad(\text{antisymmetric }\bn,\ \text{symmetric }\six).
  \label{eq:qq}
\end{equation}
A \emph{gluon} is the adjoint of $\thr\otimes\bn$:
\begin{equation}
  \thr\otimes\bn \;=\; \eg \oplus \mathbf 1 ,
  \label{eq:qqbar}
\end{equation}
which requires one quark \emph{and one antiquark}. A loop around two
\emph{quark} holes can therefore never be in the gluon octet; it is a
\textbf{diquark}. The singlet selects the antitriplet part:
$\bn\otimes\thr=\mathbf 1\oplus\eg\supset\mathbf 1$, so the diquark loop must be
the ``good'' $\bn$ to neutralize against the third quark ($\thr$). This is the
group--theoretic shadow of \S\ref{sec:duality}: the homological signature
$[\gamma_{ij}]=-[k]$ (a sign/conjugation) \emph{is} the $\bn$ being the conjugate
of the complementary $\thr$. Two derivations, Gauss law and Poincar\'e duality,
one conclusion.

\paragraph{So the within--proton loops corroborate the diquark build, not a gluon
picture.} Far from a competing reading, ``the loops are diquarks'' is the same
object as the $\bn$ diquark of the two--step build [\#489] and the diquark
sequence caveat [CW\,\S7]. The edge basis makes the diquark a \emph{first--class
geometric object} (an enclosing loop), which is conceptually cleaner than tracking
it as interior structure.

\paragraph{Where a genuine gluon loop lives.}
A gluon loop \eqref{eq:qqbar} needs a quark \emph{and} an antiquark inside it ---
and those exist in this construction. The build is a simultaneous pair--creation
event [SR\,\S7], $3$ neutral $q\bar q$ pairs $\to[\text{proton},
\text{antiproton}]$, so antiquark holes are present (each pair's $\bar q$, and the
antiproton leg). A loop enclosing a proton quark hole and a nearby antiquark hole
is in $\thr\otimes\bn=\eg\oplus\mathbf 1$: the \emph{meson/gluon channel}. So the
correct statement is:
\begin{equation}
  \boxed{\;\text{within--proton pair loop}\ \in\ \thr\otimes\thr=\bn\oplus\six\ (\text{diquark});
          \quad
          \text{quark--antiquark loop}\ \in\ \thr\otimes\bn=\eg\oplus\mathbf 1\ (\text{gluon/meson}).\;}
\end{equation}
``Loops as gluons'' is right \emph{only} for the cross--leg quark--antiquark loops,
not for the within--proton pair loops the diquark proposal uses.

\paragraph{What a gluon loop would carry --- and the constituent caveat.}
A gluon is a spin--$1$ color octet: a quark--antiquark loop carries the
\emph{chromo--flux/binding} (the $Y$--shaped baryon string between the three color
charges), contributing to mass and orbital angular momentum, not to the valence
spin--$\halff$ content. This dovetails with [CW\,\S4], which finds the entangling
exponential acting as the one--gluon--exchange \emph{hyperfine} (color--magnetic
spin--spin) interaction --- the two views are complementary: [CW] reads the gluon
as the spin--spin \emph{operator} $\bm S_i\!\cdot\!\bm S_j$ whose sign splits
$N$ from $\Delta$, while this note reads the gluon as the \emph{color rep} of a
$q\bar q$ loop. Both decline to put the proton's $\tfrac34$ on the gluon.
\emph{Aside, flagged as beyond scope:} in real QCD gluon spin and orbital angular
momentum \emph{do} carry a large part of the nucleon spin (the ``spin crisis'':
valence quarks supply only ${\sim}30\%$) \cite{spincrisis}. The $q\bar q$ loops
are exactly where such contributions would live --- but [SR]'s target is the
constituent--model $\tfrac34$ from three valence spins, and within that target the
gluon loops are the glue, not the valence spin.

% =============================================================================
\section{Caveats, reduction to the open kernel, and one cheap new test}
\label{sec:caveats}
% =============================================================================

\begin{enumerate}[leftmargin=2em]
  \item \textbf{The edge basis is not a free lunch on spin.} By \S\ref{sec:spin}
        the spin content of a pair--loop is the closed--loop holonomy / the
        correlator $C_{ij}$ [CW\,\S5--6]. Reading it still requires the
        faithful joint two--hole state, i.e.\ the same fiber$\leftrightarrow$cells
        (Whitney / K\"ahler--Atiyah) lift named open in [SR\,\S9] and [CW\,\S6].
        This note does \emph{not} bypass that kernel; it reframes the
        objects feeding it.
  \item \textbf{Period $\ne$ holonomy.} The flavor result \S\ref{sec:flavor} uses
        the \emph{linear} period; the spin result needs the \emph{bilinear}
        holonomy. Conflating them rebuilds the product floor in loop coordinates
        (the linear sum of two Bloch vectors is still a one--body object). Keep the
        two readings of a ``loop'' distinct.
  \item \textbf{Three--body, not two.} A single pair--loop is a two--body object;
        the proton is three. As in [CW\,\S7] the clean realization is the
        \emph{diquark sequence} (bond $\{1,2\}$, then couple $3$), which is what
        the two--step build already does --- and which the edge basis names as
        ``form $\gamma_{12}$, then close the remaining loops.'' There is no single
        three--body ``super--loop'' that linearizes the $\mathrm{GHZ}/W$ structure.
\end{enumerate}

\paragraph{The genuinely new, cheap experiment.}
Everything \emph{spin} here reduces to existing plans, but one prediction of the
\emph{flavor} result is new and cheap to falsify. The dual--flavor claim
\eqref{eq:loopflavor}--\eqref{eq:tau} says: read the DK taste/charge not per hole
but on the \emph{three pair--loops} (the dual basis), and you should recover
$\{d,u,u\}$ with the diquark loop $\gamma_{12}$ \emph{distinguished} as the
$\tau$--spectator. Concretely, on the built proton [\#489]:
\begin{enumerate}[leftmargin=1.9em]
  \item For each pair $\{i,j\}$, form the dual--basis period
        $w_i+w_j=\oint_{\gamma_{ij}}\psi$ from the already--computed per--hole
        weights (no new geometry), and read the DK taste on that combination.
  \item Check (a) the multiplicity is $2$:$1$ ($u$:$u$:$d$), (b) the
        \emph{odd--one--out} loop coincides with the diquark loop $\gamma_{12}$
        (equivalently is dual to the step--2 quark hole $3$), and (c) the result is
        invariant under the GAUGE/RELABEL gates [SR\,\S7]. Failure of (b) would
        falsify the dual--flavor framing immediately; success would establish that
        flavor, like the diquark build, is a Poincar\'e--dual statement on the
        register --- the one claim of this note that is both new and testable
        without the open lift.
\end{enumerate}
This is a sibling of, not a substitute for, [CW]'s pairwise--$C_{ij}$ program: that
one targets \emph{spin} (and needs the lift); this one targets \emph{flavor} (and
does not).

% =============================================================================
\section{Status and references}
\label{sec:status}
% =============================================================================

\paragraph{Status.} A lens, not a validated result. Established here: the
pair--loop$\,\to\,$complementary--hole duality \eqref{eq:dual} (homology) and the
diquark--not--gluon classification \eqref{eq:qq}--\eqref{eq:qqbar} (rep theory),
together with the consistent dual--flavor labeling \S\ref{sec:flavor}. Reframed,
not resolved: the spin readout (\S\ref{sec:spin} reduces it to [CW]). New and
testable: the dual--basis flavor read \S\ref{sec:caveats}.

\paragraph{One sentence.} ``Make the loops the quarks'' is the Poincar\'e--dual
edge basis of the same three color holes --- it carries the identical singlet,
assigns $d$ to the diquark loop precisely because that loop is dual to the
$d$--quark, and (by $\thr\otimes\thr=\bn\oplus\six$) is a diquark rather than a
gluon; it sharpens flavor and color but inherits, for spin, the same closed--loop
curvature and the same open lift the sibling notes already name.

\paragraph{Code and companions.}
\begin{itemize}[leftmargin=1.6em]
  \item Sibling documents (this directory): \code{spin\_readout.tex} [SR] (the
        study document, which names the total--space operator);
        \code{cartan\_weyl\_gluon.tex} [CW] (the Cartan/Weyl lens, $C_{ij}$, the
        hyperfine gluon, and the fiber$\leftrightarrow$cells lift with the
        closed--loop holonomy escape hatch).
  \item Modules: \code{examples/cobordism/dk\_composite\_spin.py}
        (\code{wilson\_line}, \code{facet\_transport}, the carried weights $w_h$,
        the product floor); \code{examples/cobordism/dk\_spin\_wilson.py} (the
        hinge--loop curvature); \code{dk\_joint\_spin.py},
        \code{test\_dk\_joint\_spin.py} (the validated $J^2$ instrument).
  \item Build/findings: \code{docs/design/joint\_proton\_spin\_findings.md} (\#489,
        the two--step diquark build and the validated measuring stick).
  \item Issues: \#410 (flavor/charge track); \#489 (build); \#495 / \#477 (the
        readout this companion's spin reduces to).
\end{itemize}

% =============================================================================
\begin{thebibliography}{9}
\addcontentsline{toc}{section}{References}

\bibitem{anselmino} M.~Anselmino, E.~Predazzi, S.~Ekelin, S.~Fredriksson, and
  D.~B.~Lichtenberg, \emph{Diquarks}, Rev.\ Mod.\ Phys.\ \textbf{65}, 1199 (1993)
  --- the $\bn$ ``good'' diquark and $\thr\otimes\thr=\bn\oplus\six$.

\bibitem{jaffe} R.~L.~Jaffe, \emph{Exotica}, Phys.\ Rept.\ \textbf{409}, 1 (2005)
  --- diquark correlations in baryons; the antisymmetric--antitriplet selection.

\bibitem{hatcher} A.~Hatcher, \emph{Algebraic Topology}, Cambridge Univ.\ Press
  (2002) --- homology, Poincar\'e duality, additivity of periods over cycles.

\bibitem{spincrisis} C.~A.~Aidala, S.~D.~Bass, D.~Hartline (Hatta), and
  Yu.~V.~Kovchegov, \emph{The spin structure of the nucleon}, Rev.\ Mod.\ Phys.\
  \textbf{85}, 655 (2013) --- the gluon spin / orbital angular momentum
  contributions (the ``spin crisis'').

\bibitem{siblings} Tessera theory notes (this directory):
  \emph{From Spin--\textonehalf{} to the Entangled Readout of the Proton Spin}
  (\code{spin\_readout.tex}); \emph{A Cartan / Weyl--Chamber
  Reading of the Composite--Spin Readout} (\code{cartan\_weyl\_gluon.tex}).

\end{thebibliography}

\end{document}
